\begin{abstract}
Thermodynamic experiments frequently compress an ordered constitutive process into an endpoint spectrum, a finite response table, or a small family of scalar summaries. Such a representation can be exact on the data it records while remaining blind to order information required by a downstream target. We develop a finite-dimensional thermodynamic specialization of spectral-blindness and Recognition Kernel theory.

The state of observation is a typed, nondimensionalized thermodynamic response fibre $\Rth\cong\C^m$ equipped with a declared positive response metric $\thmmetric>0$. Ordered response generators $A_j$ produce endpoint operators $U_\pi$ through labelled exponentials. We prove that ordinary characteristic determinants are exactly blind to two-factor reversal and cyclic rotation, even when the endpoint operators differ. For three labelled generators, the first multilinear connected noncyclic channel is
\[
 [xyz]\bigl(L_{CBA}-L_{CAB}\bigr)
 =\frac{q}{(1-q)^2}\Tr\bigl(A[B,C]\bigr).
\]
A declared thermodynamic marker detects an invisible endpoint residue through a nonzero metric trace pairing, and complete marker banks reconstruct endpoint differences with quantitative frame bounds.

The thermodynamic metric adds a structure absent from the bare spectral problem. For a response event $y$ and decoder $c$, the exact cut-square identity decomposes response energy into the declared scalar channel and transverse constitutive mismatch:
\[
 \|c\|_{\thmmetric}^2\|y\|_{\thmmetric}^2
 -|\langle c,y\rangle_{\thmmetric}|^2
 =\|\thmmetric^{1/2}c\wedge \thmmetric^{1/2}y\|^2.
\]
For four normalized response channels this residue is the sum of six pairwise squares. General multiplicity transforms then specialize to thermodynamic response operators: at total degree five the collapsed one-exponential-per-label model has fifty formal cyclic coordinates, functional rank forty-five, and a five-dimensional blind kernel. Five split-occurrence scalar probes repair this defect at $q=2$, and no smaller scalar bank can do so. More generally, for a declared thermodynamic target $\Pi$, the exact minimum number of supplemental scalar observations is $\rank(\Pi|_{\Ker A})$.

The results distinguish endpoint spectral blindness, target-relevant response blindness, and erased-history blindness. They do not claim that equilibrium response identities alone select noncommuting response generators, that a finite algebraic marker is automatically a physical sensor, or that endpoint observations recover histories already collapsed to the same endpoint.
\end{abstract}

\noindent\textbf{Keywords.}
thermodynamic response; spectral blindness; ordered constitutive processes; Recognition Kernel; cut-square identity; observer markers; target-faithful completion; Onsager response; minimal lifts; noisy reconstruction.

\tableofcontents
